% ID: MAT-ALG-INT-013
% Titolo: Espressione con quozienti e basi negative
% Disciplina: Matematica
% Area: Algebra
% Argomento: Numeri interi
% Classe: Prima secondaria di secondo grado
% Difficolta: 3
% Tipo: problema_numerico
% Tag: matematica, algebra, numeri_interi
% Fonte: appunti_matematica_anonimizzati
% Provenienza: materiale fornito dal docente; testo anonimizzato; soluzione redatta per Exergo

\begin{esercizio}
Calcola $\left\{(-21)^2:(-7)^2+\left[(-21)^3\right]^2:(-21)^5\right\}:(-2)^2$.
\end{esercizio}

\begin{soluzione}
Il primo quoziente è 9, il secondo è $(-21)^{6-5}=-21$. Quindi $(9-21)/4=-3$.
\end{soluzione}
